% Section II -- Related Work.
% Positioning: prior work either optimises energy under a DP constraint
% analytically, or measures one privacy setting. Neither sweeps epsilon with
% hardware counters, and none reports energy-to-target-accuracy.

\section{Related Work}

\subsection{Measuring the energy of training}

Green AI made energy and carbon reportable quantities rather than
externalities~\cite{greenai, strubell}, and a family of tools now estimates
them from utilisation traces and grid factors~\cite{henderson, carbontracker,
lacoste}. Comparisons of those tools against instrumented hardware find the
estimates diverge substantially from measurement~\cite{bouza,
aquinobritez}, which is why we read NVML counters rather than model power
from device ratings. Patterson et al.~\cite{patterson} established that
where a job runs changes its emissions by $5$--$10\times$; our geographic
result is that finding applied to a parameter, $\varepsilon$, rather than to
a schedule. Savazzi et al.~\cite{savazzi} give an analytical energy and
carbon framework for federated and decentralised learning, but without a
privacy mechanism.

\subsection{Energy and differential privacy in federated learning}

One line of work treats DP noise as a constraint inside a wireless
resource-allocation problem and minimises device energy under it, either by
harvesting channel noise as the privacy mechanism~\cite{liu, zhang, mao} or
by co-designing compression with the noise
budget~\cite{chen, ding}. Energy there is an optimisation objective evaluated
through a convergence bound; it is not measured, and $\varepsilon$ is a
constraint rather than a swept variable. Closest to us empirically,
Schwermer et al.~\cite{schwermer} built a physically distributed edge testbed
and reported traffic, energy and training time for a CNN on MNIST with
differential privacy enabled, and found software estimates deviating by up to
$35\%$ from measurement. Their contrast is DP on versus off at one setting;
ours is the shape of the cost across the budget. Aroussi
et al.~\cite{aroussi} schedule $\varepsilon$ adaptively on Raspberry Pi and
ESP32 nodes and report a $20$--$25\%$ energy improvement over a fixed-privacy
baseline --- the gain of a controller, rather than the cost curve a designer
would need to choose the budget in the first place.

\subsection{Rounds as the mechanism}

That the cost should act through the round count is known analytically. Wei
et al.~\cite{nbafl} prove that under client-level DP there is an optimal
number of aggregations for a given protection level, beyond which
convergence degrades, and later derive a discounting method to find
it~\cite{weicdp}; the same optimum reappears for personalised
FL~\cite{weipfl}, for query and reply counts~\cite{zhou}, and in control
formulations that adapt noise and rounds jointly~\cite{yuan}. These results
say where to stop. What they do not say is what continuing costs, because
the round is an abstract unit in them. Our measurements price it: at
$\varepsilon = 0.5$ on UCI-HAR, $83\%$ of a run's joules are spent after the
accuracy peak that this literature predicts.

\subsection{Our position}

We therefore contribute the missing empirical axis: hardware-counter energy
swept across $\varepsilon$ under whole-run client-level
accounting~\cite{dpsgd, geyer}, decomposed into a fixed per-round overhead
and a convergence penalty, and reported as energy-to-target-accuracy rather
than energy per round or energy per run.
